\documentclass[final]{beamer}
\usepackage[orientation=portrait,size=a1,scale=1.18]{beamerposter}
\geometry{hmargin=1.5cm,vmargin=0cm}
\usepackage[utf8]{inputenc}
\linespread{1.04}

% Custom theme
\usetheme{sharelatex}
\definecolor{primary}{RGB}{160,32,60}

% Packages
\usepackage{amsmath,amssymb}
\usepackage{graphicx}
\usepackage{tikz}
\usepackage{multicol}
\usepackage{xcolor}
\usepackage{enumitem}
\setlist[itemize]{leftmargin=1.5em, itemsep=0.35em, parsep=0pt, topsep=0.2em}
\setlength{\columnsep}{2cm}
\usepackage[backend=bibtex, style=numeric, sorting=none, maxnames=2, minnames=1]{biblatex}
\AtBeginBibliography{
  \fontsize{16.15}{7}\selectfont
  \setlength{\itemsep}{0pt}
  \setlength{\parskip}{0pt}
}
\addbibresource{bibliography.bib}

% ── Figure placeholder ────────────────────────────────────────────────────────
% \figplaceholder[height]{label}
\newcommand{\figplaceholder}[2][5.5cm]{%
  \begin{center}
    \colorbox{gray!12}{%
      \parbox[c][#1][c]{0.92\linewidth}{%
        \centering\color{gray!55}\large\itshape #2%
      }%
    }
  \end{center}
}

% ── Title block ───────────────────────────────────────────────────────────────
\title[Reinforcement Learning -- CentraleSup\'elec]{%
  Snake in All its Forms - Reinforcement Learning in Custom Snake Environments}
\author{Adonis Jamal, Fotios Kapotos, Jean-Vincent Martini}
\institute[CentraleSup\'elec]{CentraleSup\'elec}
\date{March 2026}

% =============================================================================
\begin{document}
\begin{frame}[t] 
\vskip0.1cm
\begin{multicols}{2}

% ── Block 1 : The Snake Game ──────────────────────────────────────────────────
\begin{block}{1.\enspace The Snake Game}

\begin{figure}[h]
    \centering
    \includegraphics[width=0.5\linewidth]{img/map_simple.png}
    \caption{Simple $16\!\times\!16$ grid}
    \label{fig:map_simple}
\end{figure}

\begin{itemize}
  \item Agent controls a snake on a $16\!\times\!16$ grid; eating food grows the body by one segment
  \item Terminates on wall or self-collision; Objective: Maximising the snake's length
  \item \textbf{MLP agent:} 15-dim hand-crafted input: danger flags, direction bits, food/poison bearings
  \item \textbf{CNN agent:} 6-channel raw grid (body, head, gold, silver, poison, obstacles); learns spatial patterns directly
  \item Both trained with DQN; action space: 4 absolute directions
\end{itemize}

\end{block}

\begin{figure}[h]
    \centering
    \includegraphics[width=1\linewidth]{img/train_simple.png}
    \caption{Comparaison between CNN and MLP Networks}
    \label{fig:example}
\end{figure}

% ── Block 2 : The Long Snake Problem ─────────────────────────────────────────
\begin{block}{2.\enspace The Long Snake Problem}

\figplaceholder[6cm]{Figure: snake trapping itself in a corridor, food locked inside}

\begin{itemize}
  \item More body segments $\Rightarrow$ more constraints, less free space to manoeuvre
  \item Greedy food-chasing closes off the snake's own escape routes
  \item Reward becomes sparse: death rare at length~3, near-certain at length~20+
  \item Agent must plan \textbf{many steps ahead} to keep passages open
  \item \textbf{Potential-based distance shaping:} reward proportional to the reduction in Manhattan distance to the nearest food each step
  \item \textbf{Loop penalty} ($-8$): discourages revisiting the same cell $\geq\!4$ times within a 20-step window
  \item \textbf{Curriculum:} agent trained at initial lengths $3\!\to\!7\!\to\!15$ before transferring to the hard environment
\end{itemize}

\end{block}


\columnbreak


% ── Block 3 : Hard Mode Environment ──────────────────────────────────────────
\begin{block}{3.\enspace Hard Mode: Dynamic Environment}

\begin{figure}[h]
    \centering
    \includegraphics[width=0.5\linewidth]{img/map_hard.png}
    \caption{Hard environment: gold, silver, and poison food alongside static obstacles on a $16\!\times\!16$ grid}
    \label{fig:map_hard}
\end{figure}

\begin{itemize}
  \item \textbf{Gold food} ($+5$): primary target; respawns immediately upon collection
  \item \textbf{Silver food} ($+0.5$): secondary target; worth pursuing when gold is distant
  \item \textbf{Poison food}: shrinks the snake by 2 segments upon consumption; must be actively avoided
  \item \textbf{Static obstacles}: fixed 3-cell walls spawned at episode reset; lethal on contact
  \item \textbf{Dynamic obstacles}: 3-cell walls that bounce back and forth within a fixed range; lethal even when a wall \emph{moves into} a stationary snake
  \item A single frame is insufficient to infer obstacle direction --- \textbf{temporal context is essential}
\end{itemize}

\end{block}


% ── Block 4 : Strategies ──────────────────────────────────────────────────────
\begin{alertblock}{4.\enspace Strategies for the Hard Environment}

\figplaceholder[6cm]{Figure: frame stacking ($t{-}3$ to $t$, obstacle drifting) fed into CNN; dueling DQN split ($V$ and $A$ heads)}

\begin{itemize}
  \item \textbf{Frame stacking ($n\!=\!4$):} encodes obstacle velocity; converts partial observability into an approximately Markovian state
  \item \textbf{Dueling DQN:} shared trunk splits into $V(s)$ and $A(s,a)$ heads; $Q = V + A - \bar{A}$; faster convergence when action choice is irrelevant
  \item \textbf{Double DQN:} decouples action selection from target evaluation; reduces Q-value overestimation
  \item \textbf{$n$-step returns ($n\!=\!3$):} $G_t = \sum_{k=0}^{2}\gamma^k r_{t+k}$; propagates credit further back in time
  \item \textbf{Soft target update ($\tau\!=\!0.005$):} Polyak averaging of target weights; smoother training
  \item \textbf{Distance shaping:} potential-based reward toward nearest gold/silver each step
\end{itemize}

\end{alertblock}


% ── Block 5 : Results ─────────────────────────────────────────────────────────
\begin{block}{5.\enspace Results}

\figplaceholder[8cm]{Figure: learning curves (reward vs.\ training steps) --- MLP vs.\ CNN on simple env; ablation bar chart for hard env strategies}

\begin{itemize}
  \item MLP converges faster on the simple environment; CNN generalises better spatially in the hard environment
  \item Curriculum training is essential: agents trained directly on the hard environment without warmup fail to converge
  \item Frame stacking and dueling architecture are the most impactful components for dynamic obstacle avoidance
  \item Double DQN and $n$-step returns provide consistent gains across all experimental settings
\end{itemize}

\end{block}


% ── References ────────────────────────────────────────────────────────────────
\begin{block}{References}
{\tiny\printbibliography[heading=none]}
\end{block}


\end{multicols}
\end{frame}
\end{document}
